\documentclass{EasyChair3.5/easychair}

\usepackage{microtype}
\usepackage[T1]{fontenc}
\usepackage[utf8]{inputenc}
\usepackage[noabbrev,nameinlink,capitalize,poorman]{cleveref}

\title{SAT Resetmaxxing for Fun and Profit \\ \small (Final Title TBD)}
\author{Martin S.\ Fixman\inst{1} \and Frederik Rømming\inst{1} \and Sean B.\ Holden\inst{1}}

\institute{
	\textsuperscript{1} University of Cambridge,
	United Kingdom
}

\titlerunning{SAT Resetmaxxing for Fun and Profix}
\authorrunning{Fixman, Rømming, Holden}

\bibliographystyle{plain}

\newcommand{\todo}[1]{\textcolor{red}{(TODO: #1)}}
\newcommand{\step}[1]{\textbf{#1}}

\begin{document}
\maketitle{}

\abstract{%
	We improve SATCoP performance by resetting the connection tableau often.
	This causes more and more varied SAT statements to be added to the SAT prover, which causes it to saturate faster.
}

\section{Introduction}

\section{Preliminaries}

\todo{Tableau or Tableaux? Bibliography is inconsistent.}

\subsection{Tableau Solving Methods}

Minimalist tableau-style provers based on Prolog have been used with great success since the release of leanTAP in 1995\cite{beckertLeanTAPLeanTableaubased1995}, and have consistently achieved high rankings in solving FOF and CNF problems in the Thousands of Problems for Theorem Provers (TPTP) library\cite{sutcliffe12thIJCARAutomated2025}.

Rather than working like saturation provers \todo{what do saturation provers even \emph{do}?} such as Vampire prover\cite{riazanovDesignImplementationVAMPIRE2002}, Tableau-based provers attempt to create a Connection Tableau with ground terms.
The prover builds a tree, where each node corresponds to a particular clause of the problem in conjunctive normal form, using one of the following four steps.


\begin{figure}[h]
	\fbox{\parbox[c][75pt]{\textwidth}{\centering \todo{Add a diagram and/or a description.}}}
	\caption{Description of the operations of Tableau calculus\cite{farberCuriouslyEffectiveBacktracking2024}.}
\end{figure}

LeanCoP\cite{ottenLeanCoPLeanConnectionbased2003} improves 

Talk about LeanCoP\cite{ottenLeanCoP20}, possibly MeanCoP\cite{farberCuriouslyEffectiveBacktracking2024}.

\subsection{SATCoP and Similar Methods}

\section{Research Vehicle}

Talk about the Connections Prover\cite{connections_2023}.

\section{Experimental Evaluation}

\section{Future Results}

\bibliography{bibliography}

\end{document}
